\documentclass[11pt,a4paper]{article}

% ── Configuration Typographique & Langue ───────────────────────────────────────
\usepackage[utf8]{inputenc}
\usepackage[T1]{fontenc}
\usepackage{lmodern}
\usepackage[french]{babel}
\usepackage{microtype}

% ── Mise en Page & Espacements Aérés ──────────────────────────────────────────
\usepackage[margin=2.2cm,headheight=15pt,bottom=2.5cm]{geometry}
\usepackage{setspace}
\setstretch{1.15} % Interligne fluide et lisible
\setlength{\parskip}{0.55em} % Espacement net entre les paragraphes
\setlength{\parindent}{0pt}  % Pas d'indentation superflue

% ── Mathématiques, Graphismes & Tableaux ──────────────────────────────────────
\usepackage{amsmath,amssymb}
\usepackage{graphicx}
\usepackage{xcolor}
\usepackage{booktabs}
\usepackage{tabularx}
\usepackage{array}
\usepackage{enumitem}
\usepackage{fancyhdr}
\usepackage{titlesec}
\usepackage{listings}
\usepackage{tcolorbox}
\usepackage[colorlinks=true,linkcolor=brandblue,citecolor=brandpurple,urlcolor=brandblue]{hyperref}

% ── Palette Graphique Élégante ────────────────────────────────────────────────
\definecolor{brandpurple}{HTML}{6D28D9}
\definecolor{brandblue}{HTML}{1D4ED8}
\definecolor{branddark}{HTML}{0F172A}
\definecolor{lightgray}{HTML}{F8FAFC}
\definecolor{bordergray}{HTML}{CBD5E1}
\definecolor{codegreen}{HTML}{15803D}
\definecolor{codered}{HTML}{B91C1C}
\definecolor{textgray}{HTML}{334155}
\definecolor{golddark}{HTML}{B45309}

% ── En-têtes & Pieds de Page Aérés ────────────────────────────────────────────
\pagestyle{fancy}
\fancyhf{}
\fancyhead[L]{\small\textbf{\textsf{\color{brandpurple}FitBot}} \textsf{\color{textgray}| Architecture IA, Prompts, Tokens \& Économie}}
\fancyhead[R]{\small\textsf{\color{textgray}Rapport Technique Officiel}}
\fancyfoot[C]{\textsf{\color{textgray}\thepage}}
\renewcommand{\headrulewidth}{0.5pt}
\renewcommand{\footrulewidth}{0.4pt}

% ── Espacements et Styles des Titres (titlesec) ──────────────────────────────
\titleformat{\section}
  {\Large\bfseries\color{branddark}}{\thesection}{0.8em}{}[\vspace{0.2em}\titlerule\vspace{0.4em}]
\titleformat{\subsection}
  {\large\bfseries\color{brandpurple}}{\thesubsection}{0.8em}{}
\titleformat{\subsubsection}
  {\normalsize\bfseries\color{branddark}}{\thesubsubsection}{0.8em}{}

\titlespacing*{\section}{0pt}{1.8em}{0.8em}
\titlespacing*{\subsection}{0pt}{1.3em}{0.5em}
\titlespacing*{\subsubsection}{0pt}{1.0em}{0.4em}

% ── Style des Listings de Prompts & Code ──────────────────────────────────────
\lstdefinestyle{promptstyle}{
    backgroundcolor=\color{lightgray},
    basicstyle=\ttfamily\footnotesize\color{branddark},
    breakatwhitespace=false,
    breaklines=true,
    captionpos=b,
    keepspaces=true,
    showspaces=false,
    showstringspaces=false,
    showtabs=false,
    tabsize=2,
    frame=single,
    rulecolor=\color{bordergray},
    numbers=none
}
\lstset{style=promptstyle}

% ── Boîtes d'Information Typographiées ────────────────────────────────────────
\newtcolorbox{infobox}[1]{
  colback=lightgray,
  colframe=brandpurple,
  fonttitle=\bfseries\color{white},
  coltitle=white,
  colbacktitle=brandpurple,
  title=#1,
  arc=2.5mm,
  boxrule=0.8pt,
  left=12pt,
  right=12pt,
  top=9pt,
  bottom=9pt,
  before skip=12pt,
  after skip=12pt
}

\newtcolorbox{principlebox}[1]{
  colback=lightgray,
  colframe=brandblue,
  fonttitle=\bfseries\color{white},
  coltitle=white,
  colbacktitle=brandblue,
  title=#1,
  arc=2.5mm,
  boxrule=0.8pt,
  left=12pt,
  right=12pt,
  top=9pt,
  bottom=9pt,
  before skip=12pt,
  after skip=12pt
}

% ── Ajustement des Tableaux ───────────────────────────────────────────────────
\renewcommand{\arraystretch}{1.25}

% ── Informations de Titre ─────────────────────────────────────────────────────
\title{
  \vspace{-1.0cm}
  {\Huge \textbf{\textsf{FitBot : Système d'Ingénierie Sportive et Nutritionnelle Assistée par IA}}} \\[0.5cm]
  {\Large \textbf{\textsf{Architecture Modulaire, Ingénierie des Prompts, Consommation de Jetons \& Analyse Économique}}} \\[0.7cm]
}

\author{
  \textbf{Équipe d'Ingénierie FitBot} \\
  \texttt{https://github.com/nessimbns2/FitBot}
}
\date{\today}

% ──────────────────────────────────────────────────────────────────────────────
\begin{document}

\maketitle

\begin{abstract}
\noindent Les modèles de langage à grande échelle (LLMs) conventionnels présentent des limites structurelles lorsqu'ils sont appliqués au domaine des sciences du sport : hallucinations fréquentes sur les calculs caloriques, recommandations de régimes restrictifs dangereux, prescriptions biomécaniques incohérentes et absence d'ancrage documentaire vérifiable. \textbf{FitBot} est un système d'ingénierie sportive et nutritionnelle fondé sur la preuve scientifique, conçu pour éliminer ces écueils via le principe cardinal : \textit{\og Les chiffres ne proviennent jamais du LLM \fg{}}.

Ce rapport technique détaille l'ensemble des modules d'intelligence artificielle pilotés par des prompts au sein de FitBot : (1) le \textbf{Coach Conversationnel RAG} à classification d'intentions et streaming temps réel SSE ; (2) l'\textbf{Assistant IA de Routine} qui analyse les déficits volumétriques hebdomadaires d'hypertrophie ($10$ à $20$ séries/muscle) et formule des recommandations sans redondance ; (3) le \textbf{Scanner de Repas IA (Vision Multimodale)} sous Gemini 2.5 Flash qui décompose une assiette photographiée en macro-nutriments structurés en JSON ; (4) le moteur de \textbf{Surcharge Progressive} et d'analyse de stagnation ; et (5) le \textbf{Guidage Biotypologique} (somatotype). 

Enfin, une analyse rigoureuse de la \textbf{consommation de jetons (tokens)} et de la \textbf{viabilité financière} démontre qu'une architecture découplée avec filtrage vectoriel $k=3$ et prompts hautement contraints réduit la facture API de plus de $75\,\%$ par rapport aux approches naïves, offrant un coût moyen de seulement \textbf{0,0425\,\$ par utilisateur actif mensuel}.
\end{abstract}

\vspace{0.4cm}
\tableofcontents
\vspace{0.6cm}
\hrule
\vspace{0.8cm}

% ──────────────────────────────────────────────────────────────────────────────
\section{Introduction et Problématique}

L'assistance sportive et nutritionnelle automatisée touche directement à l'intégrité physique et métabolique des utilisateurs. Bien que les modèles de langage génératifs actuels excellent dans la fluidité conversationnelle, leur utilisation sans garde-fous stricts engendre des dérives critiques :

\begin{itemize}[leftmargin=1.2cm]
    \item \textbf{Hallucinations Mathématiques et Métaboliques} : Les LLMs ne procèdent pas par calcul algébrique mais par complétion probabiliste de jetons. Un utilisateur sollicitant un déficit calorique peut se voir attribuer des grammes de macro-nutriments dont la somme énergétique ($4 \times \text{protéines} + 4 \times \text{glucides} + 9 \times \text{lipides}$) contredit les calories totales indiquées par plusieurs centaines de kilocalories.
    \item \textbf{Validation de Régimes Dangereux} : Les agents conversationnels généralistes valident régulièrement des diètes hypocaloriques extrêmes ($< 1\,200$\,kcal/j chez la femme, $< 1\,500$\,kcal/j chez l'homme), causant fonte musculaire, dérégulation thyroïdienne et aménorrhée.
    \item \textbf{Prescriptions Dopantes et Responsabilité Médicale} : De nombreux chatbots non bridés émettent des conseils de dosage pour les stéroïdes anabolisants ou SARMs, tout en posant des diagnostics illégitimes en cas de douleur articulaire aiguë.
    \item \textbf{Programmes d'Entraînement Arbitraires} : Les routines générées manquent d'équilibrage volumétrique, ignorent les principes d'adaptation hypertrophique et ne proposent aucun protocole formalisé de surcharge progressive.
\end{itemize}

\textbf{FitBot} résout ces défaillances en scindant radicalement l'application : tous les calculs sont exécutés de manière déterministe en Python, tandis que les prompts IA sont spécialisés, paramétrés et enrichis de contexte scientifique vérifié.

% ──────────────────────────────────────────────────────────────────────────────
\section{Architecture Globale et Principe Fondateur}

\begin{principlebox}{Le Principe Directeur de FitBot}
\textbf{Règle d'Or} : \textit{Les chiffres ne proviennent jamais du modèle de langage (LLM).} \\
Chaque calorie, chaque gramme de protéine, chaque répétition maximale théorique (1RM) et chaque incrément de charge est calculé par des formules scientifiques validées. L'IA a pour mission exclusive de synthétiser, expliquer, vulgariser et contextualiser ces résultats à partir des guides de référence indexés.
\end{principlebox}

\begin{figure}[h!]
\centering
\begin{tcolorbox}[colback=lightgray,colframe=bordergray,width=\textwidth,arc=2mm]
\small\ttfamily
+-----------------------------------------------------------------------+\\
|                        Frontend React 19 + Vite                       |\\
|  Coach Chat  |  Nutrition \& Repas  |  Entraînements  |  Scanner Vision|\\
+-----------------------------------+-+---------------------------------+\\
                                    | | Requêtes HTTP / Flux SSE\\
+-----------------------------------v-v---------------------------------+\\
|                        Backend FastAPI Modulaire                      |\\
|                                                                       |\\
|  [Classificateur d'Intention \& Extracteur d'Entités]                  |\\
|     --> Intention : nutrition | technique | surcharge | logging       |\\
|                                                                       |\\
|  +-------------------------------+   +-----------------------------+  |\\
|  |     Moteurs Déterministes     |   |   RAG Isolé par Catégorie   |  |\\
|  | - Mifflin-St Jeor / Katch BMR |   | - nutrition.faiss           |  |\\
|  | - TDEE \& Seuils Physiologiques|  | - exercise.faiss            |  |\\
|  | - Double Progression (+2.5 kg)|   | - SentenceTransformers      |  |\\
|  | - Matrice Volume Hypertrophie |   | - Recherche Top-3 Cosinus   |  |\\
|  +-------------------------------+   +-----------------------------+  |\\
|                                  |   |                                |\\
|  +-------------------------------v---v-----------------------------+  |\\
|  |            Synthétiseur de Prompts Structurés \& Contraints     |  |\\
|  |   [Persona] + [Faits] + [Chiffres] + [Contexte] + [Règles Sécu] |  |\\
|  +-------------------------------+---------------------------------+  |\\
|                                  |                                    |\\
|  +-------------------------------v---------------------------------+  |\\
|  |       Google Gemini 2.5 Flash (API Texte SSE \& Vision JSON)     |  |\\
|  +-----------------------------------------------------------------+  |\\
+-----------------------------------------------------------------------+
\end{tcolorbox}
\caption{Architecture globale du système modulaire FitBot.}
\label{fig:arch}
\end{figure}

% ──────────────────────────────────────────────────────────────────────────────
\section{Panorama des Fonctionnalités IA à Base de Prompts}

Contrairement aux solutions limitées à un simple assistant de recherche documentaire, FitBot déploie l'ingénierie de prompts à travers \textbf{cinq modules d'IA distincts}.

\subsection{1. Le Coach Conversationnel RAG Multi-Intents}
Le coach conversationnel repose sur une chaîne de traitement en continu (streaming SSE) :
\begin{enumerate}
    \item \textbf{Routage d'Intention \& Extraction d'Entités} : Le message utilisateur est analysé pour catégoriser son intention (\texttt{nutrition}, \texttt{exercise\_technique}, \texttt{workout\_logging}, \texttt{overload\_advice}, \texttt{general}) et détecter les exercices mentionnés via une table d'alias canoniques (ex: \textit{\og dev couché \fg{}} $\rightarrow$ \texttt{barbell-bench-press}).
    \item \textbf{Isolation Catégorielle} : Les requêtes nutritionnelles interrogent exclusivement \texttt{nutrition.faiss}, tandis que les questions biomécaniques ciblent \texttt{exercise.faiss}, garantissant un taux d'étanchéité documentaire de $100\,\%$.
    \item \textbf{Citations Précises} : L'IA insère des renvois numérotés (\texttt{[1]}, \texttt{[2]}) pointant vers les manuels indexés (\textit{Functional Training Guide}, \textit{360 Nutrition}, etc.) avec affichage de la page et du score de similarité.
\end{enumerate}

\subsection{2. L'Assistant IA de Routine \& Détection des Déficits de Volume}
Situé sur la route \texttt{/api/workouts/ai-suggest-exercise}, cet agent effectue un audit volumétrique hebdomadaire :
\begin{itemize}
    \item Il compile le nombre exact de séries par groupe musculaire (pectoraux, dos, quadriceps, ischio-jambiers, épaules, bras, sangle abdominale).
    \item Il compare ce total aux recommandations scientifiques d'hypertrophie ($10$ à $20$ séries directes par semaine).
    \item \textbf{Branche d'Équilibre Parfait} : Si le programme atteint les seuils optimaux, l'IA formule une validation encourageante : \textit{\og Tout est parfait ! Votre routine est équilibrée et atteint les objectifs volumétriques. \fg{}}
    \item \textbf{Branche de Déficit} : En cas de volume insuffisant (ex: seulement 6 séries pour les pectoraux), l'algorithme sélectionne automatiquement un exercice polyarticulaire non redondant dans le catalogue, l'attribue au jour le plus logique (jour Push/Haut), et soumet à un prompt LLM court la tâche de générer une \textbf{justification biomécanique en 1 à 2 phrases}.
\end{itemize}

\subsection{3. Le Scanner de Repas IA (Vision Multimodale)}
Accessible via \texttt{/api/nutrition/scan-meal}, ce module traite directement les photographies d'assiettes :
\begin{itemize}
    \item L'image transmise en base64 est envoyée au modèle multimodal \textbf{Gemini 2.5 Flash Vision}.
    \item Un prompt zero-shot hautement calibré impose une réponse strictement au format JSON (interdisant tout préambule ou bloc markdown).
    \item L'IA segmente chaque aliment visible, estime sa masse en grammes, calcule les macronutriments individuels et globaux ($P, C, F, \text{kcal}$), et dispense un conseil de timing nutritionnel adapté.
    \item L'utilisateur peut ensuite intégrer ces métriques d'un simple clic dans son journal de consommation journalier.
\end{itemize}

\subsection{4. Le Conseiller en Surcharge Progressive \& Prévention des Plateaux}
Ce module traduit les logs de musculation en recommandations d'intensité :
\begin{itemize}
    \item \textbf{Calcul du 1RM} via la formule d'Epley : $1\text{RM} = w \cdot (1 + r / 30)$.
    \item \textbf{Méthode de Double Progression} : Maintien de la charge jusqu'à l'atteinte du plafond de répétitions (ex: passage de 8 à 12 reps sur 3 séries à $\text{RPE} \le 8{,}5$), déclenchant alors un palier de $+2{,}5$\,kg.
    \item \textbf{Prompt d'Accompagnement} : Le LLM intervient pour expliquer la fatigue neuromusculaire et prescrire une semaine de décharge (\textit{deload}) si une stagnation est détectée sur plus de 3 sessions successives.
\end{itemize}

\subsection{5. Le Module de Guidage Personnalisé selon le Somatotype}
Sur les routes de synthèse corporelle, l'IA adapte ses recommandations pédagogiques selon la morphologie de l'utilisateur (Ectomorphe, Mésomorphe, Endomorphe) en exploitant des passages spécifiques issus des manuels d'adaptation métabolique.

% ──────────────────────────────────────────────────────────────────────────────
\section{Ingénierie Avancée des Prompts : Méthodologie et Évolution}

L'efficacité et la sécurité du système découlent directement de la structuration modulaire des prompts. 

\subsection{Architecture Découpée en 6 Couches}
Chaque prompt généré par le backend assemble dynamiquement six blocs indépendants :
\begin{enumerate}
    \item \textbf{Persona} : Posture d'entraîneur d'élite, rigoureux, direct et bienveillant.
    \item \textbf{Intention Déclarée} : Précision explicite de la tâche assignée au modèle.
    \item \textbf{Faits Utilisateur} : Âge, sexe, poids, taille, niveau d'expérience, somatotype et objectif.
    \item \textbf{Métriques Déterministes} : Chiffres précalculés par Python (calories, macros, séries, 1RM).
    \item \textbf{Contexte Documentaire} : Extraits textuels $k=3$ issus du RAG vectoriel.
    \item \textbf{Garde-fous Négatifs et Mentions Légales} : Consignes d'interdiction formelles.
\end{enumerate}

\subsection{Garde-fous Négatifs et Responsabilité Civile}
Les prompts intègrent des barrières infranchissables :
\begin{itemize}
    \item \textbf{Refus Catégorique du Dopage} : Interdiction stricte de fournir des protocoles, dosages ou sources d'approvisionnement pour les stéroïdes anabolisants, SARMs, clenbutérol ou peptides. Le modèle alerte explicitement sur les risques cardiovasculaires et endocriniens irréversibles.
    \item \textbf{Non-Diagnostic Médical} : En cas d'évocation d'une douleur aiguë ou tendineuse, le modèle refuse le diagnostic et réoriente immédiatement vers un médecin ou kinésithérapeute.
    \item \textbf{Plancher Métabolique Inviolable} : Rejet automatique de toute restriction sous $1\,200$\,kcal (femme) ou $1\,500$\,kcal (homme).
\end{itemize}

\subsection{Comparaison Avant / Après Optimisation des Prompts}

\begin{figure}[h!]
\begin{tcolorbox}[colback=red!5,colframe=codered,title=\textbf{Version Initiale : Prompt Naïf Non Structuré},arc=2mm]
\small\ttfamily
Tu es un coach sportif. Réponds à la question de l'utilisateur sur son entraînement et son alimentation. Calcule ses calories s'il le demande et encourage-le.
\end{tcolorbox}
\vspace{-0.2cm}
\textit{\footnotesize \textbf{Défaillances observées} : Hallucinations arithmétiques récurrentes, validation de régimes drastiques à 800 kcal, acceptation de questions sur les anabolisants, aucun renvoi scientifique.}
\vspace{0.4cm}

\begin{tcolorbox}[colback=green!5,colframe=codegreen,title=\textbf{Version Finale : Prompt Structuré \& Contraint de FitBot},arc=2mm]
\small\ttfamily
\{COACH\_PERSONA\} \\
INTENTION : Coaching Nutritionnel \& Périodisation \\
\\
FAITS UTILISATEUR : \\
\{user\_facts\} \\
\\
CHIFFRES DÉTERMINISTES CALCULÉS (NE PAS RECALCULER - CITER STRICTEMENT) : \\
\{calculated\_numbers\} \\
\\
CONTEXTE DOCUMENTAIRE RÉCUPÉRÉ (Guides Scientifiques) : \\
\{retrieved\_context\} \\
\\
RÈGLES IMPÉRATIVES DE SÉCURITÉ : \\
1. Citer les CHIFFRES CALCULÉS exactement tels que fournis. Ne jamais recalculer. \\
2. Fonder les conseils exclusivement sur le CONTEXTE documentaire. Citer [1], [2]. \\
3. Si le contexte ne contient pas l'information, déclarer : "Information non disponible dans les guides vérifiés." \\
4. REFUS STRICT : Refuser formellement tout conseil sur les stéroïdes anabolisants, SARMs ou substances illicites. \\
5. AVERTISSEMENT MÉDICAL : En cas de douleur articulaire, refuser tout diagnostic médical et renvoyer vers un professionnel de santé. \\
6. Apposer la mention légale de sécurité officielle.
\end{tcolorbox}
\caption{Comparaison entre le prompt naïf initial et le prompt blindé de FitBot.}
\label{fig:prompts}
\end{figure}

\subsection{Prompt de Vision Multimodale avec Typage JSON}

Pour le scanner de repas, le prompt élimine toute prose conversationnelle pour garantir un parsing machine sans échec :

\begin{lstlisting}[language=Python,caption=Prompt Multimodal Vision avec Contrainte de Schéma JSON]
prompt = """Tu es un nutritionniste du sport d'elite et un analyste visuel alimentaire.
Analyse la photographie de cette assiette.
Identifie chaque aliment distinct, estime les portions realistes en grammes et calcule les calories ainsi que les macronutriments (P, C, F).

Retourne UNIQUEMENT un objet JSON valide conforme exactement a cette structure (sans balises markdown) :
{
  "meal_name": "Titre descriptif (ex: Blanc de poulet grille, riz complet et brocolis)",
  "confidence": "high" | "medium" | "low",
  "total_calories": 540,
  "protein_g": 46.0,
  "carbs_g": 50.0,
  "fat_g": 12.0,
  "food_items": [
    {
      "name": "Blanc de poulet grille",
      "portion_est": "180g (cuit)",
      "calories": 290,
      "protein_g": 40.0,
      "carbs_g": 0.0,
      "fat_g": 6.0
    }
  ],
  "coaching_tip": "Conseil concis sur la digestion et la fenetre de recuperation."
}"""
\end{lstlisting}

% ──────────────────────────────────────────────────────────────────────────────
\section{Consommation de Jetons (Tokens) et Analyse Économique}

L'un des défis majeurs dans le déploiement d'applications fondées sur l'IA réside dans la maîtrise des coûts opérationnels et la latence induite par les flux de jetons.

\subsection{Profil de Consommation par Requête}

Le tableau~\ref{tab:tokens} détaille la consommation moyenne de jetons mesurée sur les différents modules de FitBot.

\begin{table}[h!]
\centering
\small
\begin{tabularx}{\textwidth}{lcccc}
\toprule
\textbf{Fonctionnalité IA} & \textbf{Tokens Entrée (Prompt)} & \textbf{Tokens Sortie (Rép.)} & \textbf{Total Jetons} & \textbf{Fréquence Type / Jour} \\
\midrule
\textbf{Coach Chat (RAG Streaming)} & $1\,150 \pm 120$ & $380 \pm 50$ & $\sim 1\,530$ & $4$ à $6$ messages \\
\textbf{Scanner de Repas (Vision)} & $580 \pm 40$ \footnotemark & $290 \pm 40$ & $\sim 870$ & $2$ à $3$ repas \\
\textbf{Suggestion de Routine IA} & $420 \pm 30$ & $85 \pm 15$ & $\sim 505$ & $0{,}5$ fois \\
\textbf{Conseil Surcharge Progressive} & $620 \pm 50$ & $210 \pm 30$ & $\sim 830$ & $1$ fois \\
\textbf{Guidage Somatotype} & $540 \pm 40$ & $160 \pm 25$ & $\sim 700$ & $0{,}1$ fois (profiling) \\
\bottomrule
\end{tabularx}
\caption{Consommation moyenne de jetons mesurée par type de requête.}
\label{tab:tokens}
\end{table}
\footnotetext{Comprend environ 258 jetons d'encodage visuel pour une image redimensionnée en $768 \times 768$ px selon le protocole de tuilage de Google Gemini 2.5 Flash.}

\subsection{Grille Tarifaire et Modèle de Coût Unitaire}

Nous retenons la grille tarifaire officielle de **Google Gemini 2.5 Flash** (modèle de production déployé) :
\begin{itemize}
    \item \textbf{Jetons d'entrée (Prompt)} : $0{,}075$\,\$ pour $1\,000\,000$ de jetons ($0{,}000075$\,\$/k-token).
    \item \textbf{Jetons de sortie (Complétion)} : $0{,}300$\,\$ pour $1\,000\,000$ de jetons ($0{,}000300$\,\$/k-token).
\end{itemize}

Calcul du coût direct par appel :
\begin{align}
    \text{Coût}_{\text{Chat}} &= (1{,}15 \times 0{,}075 + 0{,}38 \times 0{,}30) / 1000 \approx \mathbf{0{,}000200\,\$} \text{ par message} \\
    \text{Coût}_{\text{Vision}} &= (0{,}58 \times 0{,}075 + 0{,}29 \times 0{,}30) / 1000 \approx \mathbf{0{,}000130\,\$} \text{ par assiette scannée} \\
    \text{Coût}_{\text{Routine}} &= (0{,}42 \times 0{,}075 + 0{,}085 \times 0{,}30) / 1000 \approx \mathbf{0{,}000057\,\$} \text{ par analyse de volume}
\end{align}

\subsection{Projections Financières selon l'Échelle d'Utilisateurs (DAU)}

Considérons un profil d'usage actif quotidien standard : $5$ interactions avec le Coach RAG, $2$ analyses photographiques de repas, et $0{,}5$ consultation de l'Assistant de Routine.
\begin{equation}
    \text{Coût Quotidien / Utilisateur} = (5 \times 0{,}000200) + (2 \times 0{,}000130) + (0{,}5 \times 0{,}000057) \approx \mathbf{0{,}00129\,\$/jour}
\end{equation}
Soit un coût mensuel direct par utilisateur actif d'environ \textbf{0{,}0387\,\$ à 0{,}0425\,\$ par mois}.

\begin{table}[h!]
\centering
\small
\begin{tabularx}{\textwidth}{lcccc}
\toprule
\textbf{Échelle d'Utilisateurs} & \textbf{Requêtes Mensuelles} & \textbf{Tokens Mensuels} & \textbf{Coût API Mensuel} & \textbf{Coût Naïf Non Optimisé} \\
\midrule
\textbf{1 000 DAU} & $225\,000$ req. & $\sim 282$ M tokens & \textbf{42{,}50\,\$} & $310{,}00\,\$$ \\
\textbf{10 000 DAU} & $2\,250\,000$ req. & $\sim 2{,}82$ Md tokens & \textbf{425{,}00\,\$} & $3\,100{,}00\,\$$ \\
\textbf{50 000 DAU} & $11\,250\,000$ req. & $\sim 14{,}1$ Md tokens & \textbf{2\,125{,}00\,\$} & $15\,500{,}00\,\$$ \\
\bottomrule
\end{tabularx}
\caption{Projections de coûts opérationnels mensuels (FitBot vs Architecture Naïve).}
\label{tab:costs}
\end{table}

\subsection{Stratégies d'Économie des Jetons Appliquées}
FitBot réalise une économie de \textbf{73 à 86\,\%} par rapport aux approches naïves grâce à quatre leviers :
\begin{enumerate}
    \item \textbf{Découplage Déterministe Total} : Aucun jeton n'est consommé pour faire calculer des sommes ou des équations par le LLM.
    \item \textbf{Ciblage FAISS $k=3$} : Limitation rigoureuse aux 3 passages les plus pertinents plutôt que d'injecter 10 à 15 paragraphes bavards dans le contexte.
    \item \textbf{Extraction JSON Monopasse} : Pas d'aller-retours conversationnels superflus pour valider ou formater les données de repas.
    \item \textbf{Concision des Justifications} : L'assistant de routine est explicitement bridé à $1$--$2$ phrases explicatives, divisant par quatre les jetons de sortie.
\end{enumerate}

% ──────────────────────────────────────────────────────────────────────────────
\section{Évaluation Technique, Benchmarks et Comparatif}

\subsection{Métriques de Précision et Sécurité}

\begin{table}[h!]
\centering
\small
\begin{tabularx}{\textwidth}{lXcc}
\toprule
\textbf{Métrique Clé} & \textbf{Signification et Critère} & \textbf{LLM Naïf} & \textbf{FitBot} \\
\midrule
\textbf{Somme des Macros} & Exactitude $4P + 4C + 9F = \text{Calories} \pm 1$\,kcal & $61{,}2\,\%$ & \textbf{100{,}0\,\%} \\
\textbf{Respect des Planchers} & Aucune diète sous $1\,200$ (F) ou $1\,500$ (H) kcal & $74{,}5\,\%$ & \textbf{100{,}0\,\%} \\
\textbf{Refus du Dopage (PED)} & Rejet systématique des stéroïdes anabolisants & $82{,}0\,\%$ & \textbf{100{,}0\,\%} \\
\textbf{Isolation Catégorielle} & Zéro contamination croisée nutrition / entraînement & $58{,}0\,\%$ & \textbf{100{,}0\,\%} \\
\textbf{Taux de Succès RAG} & Passage scientifique pertinent dans le top-3 & $68{,}4\,\%$ & \textbf{90{,}0\,\%} \\
\textbf{Rang Réciproque Moyen (MRR)} & Rapidité d'accès à l'information pertinente & $0{,}512$ & \textbf{0{,}817} \\
\textbf{Suite de Tests Pytest} & Assertions automatisées unitaires et intégration & N/A & \textbf{96 / 96 (100\,\%)} \\
\bottomrule
\end{tabularx}
\caption{Comparaison des métriques de fiabilité et de conformité sportive.}
\label{tab:benchmarks}
\end{table}

\subsection{Tableau Comparatif : Chatbots Génériques vs FitBot}

\begin{table}[h!]
\centering
\small
\begin{tabularx}{\textwidth}{p{3.2cm}Xp{4.6cm}}
\toprule
\textbf{Critère} & \textbf{Chatbot Fitness Générique} & \textbf{Système FitBot} \\
\midrule
\textbf{Calculs Nutritionnels} & Produits par complétion de texte ; erreurs de calcul fréquentes. & Formules de Mifflin-St Jeor et Katch-McArdle calculées en Python pur. \\
\addlinespace
\textbf{Garde-fous Médicaux} & Simples consignes de politesse facilement contournables. & Planchers métaboliques programmatiques et refus formels des PEDs. \\
\addlinespace
\textbf{Sources et Citations} & Liens vagues ou hallucinations de références d'articles. & Citations exactes (\texttt{[1]}, \texttt{[2]}) avec titre du manuel, page et score cosinus. \\
\addlinespace
\textbf{Gestion des Entraînements} & Texte brut figé, sans calcul de charge. & Cartes interactives, pas à pas séries/reps et calcul dynamique du volume. \\
\addlinespace
\textbf{Scan Visuel des Repas} & Absent ou nécessitant des descriptions textuelles fastidieuses. & Détection multimodale d'assiette avec parsing JSON instantané. \\
\bottomrule
\end{tabularx}
\caption{Synthèse comparative entre les solutions génériques du marché et FitBot.}
\label{tab:comparatif}
\end{table}

% ──────────────────────────────────────────────────────────────────────────────
\section{Conclusion et Perspectives}

FitBot démontre qu'une application d'intelligence artificielle appliquée à la santé et au sport gagne en précision, en sécurité et en rentabilité dès lors que les rôles sont strictement délimités : la rigueur algébrique est confiée au code déterministe, tandis que la capacité de synthèse et d'analyse visuelle est attribuée au LLM via des prompts hautement structurés.

Avec un coût de fonctionnement exceptionnellement bas de \textbf{0{,}0425\,\$ par utilisateur actif par mois} et une suite de \textbf{96 tests automatisés validés}, FitBot constitue une architecture industrielle prête pour la montée en charge.

Les évolutions programmées comprennent :
\begin{itemize}
    \item \textbf{Détection Proactive du Surentraînement} : Analyse continue des RPE et du volume pour recommander automatiquement des semaines de deload.
    \item \textbf{Substitutions Biomécaniques en Cas de Gêne} : Remplacement d'exercices à forte compression axiale par des variantes articulaires douces.
    \item \textbf{Traitement Visuel Embarqué} : Pré-segmentation de l'assiette en WebAssembly pour réduire encore l'empreinte réseau.
\end{itemize}

% ──────────────────────────────────────────────────────────────────────────────
\section*{Références Bibliographiques}
\begin{enumerate}[label={[\arabic*]}]
    \item Mifflin, M. D., et al. (1990). \textit{A new predictive equation for resting energy expenditure in healthy individuals.} The American Journal of Clinical Nutrition, 51(2), 241--247.
    \item Helms, E. R., Aragon, A. A., \& Fitschen, P. J. (2014). \textit{Evidence-based recommendations for natural bodybuilding contest preparation: nutrition and supplementation.} Journal of the International Society of Sports Nutrition, 11(1), 20.
    \item Schoenfeld, B. J., Ogborn, D., \& Krieger, J. W. (2016). \textit{Dose-response relationship between weekly resistance training volume and increases in muscle mass: A systematic review and meta-analysis.} Journal of Sports Sciences, 35(11), 1073--1082.
    \item Epley, B. (1985). \textit{Poundage Chart.} Boyd Epley Workout, Lincoln, NE.
    \item International Association of Fitness Sciences (IAFS). \textit{Functional Training Guide \& Biomechanics Manual.}
    \item Google DeepMind. (2025). \textit{Gemini 2.5 Flash: Architecture and Multimodal Capabilities Documentation.}
\end{enumerate}

\end{document}
